\documentclass[11pt]{scrartcl}

\usepackage[stixtwo]{fontsetup}
\usepackage{amsmath}
\usepackage{xfrac}

\begin{document}
A számláló és a nevező gyöktényezős alakban:
\[
	x^2 - 9 = (x - 3)(x + 3),
	\qquad
	x^2 - x - 6 = (x - 3)(x + 2),
\]
így $x \notin \{-2; 3\}$ esetén
\[
	f(x) = \frac{x + 3}{x + 2}
	\text.
\]
\begin{enumerate}
	\item $\displaystyle\lim_{x \to x_0} f(x) = A$, ha
	      $\forall \varepsilon > 0$ esetén $\exists \delta(\varepsilon) > 0$,
	      hogy $|f(x) - A| < \varepsilon$, ha $x \in D_f$ és
	      $0 < |x - x_0| < \delta(\varepsilon)$.

	      Helyettesítsünk be a definícióban szereplő egyenlőtlenségbe:
	      \[
		      \left| \frac{x + 3}{x + 2} - \frac{3}{2} \right|
		      = \left| \frac{2x + 6 - 3x - 6}{2 (x + 2)} \right|
		      = \frac{|x|}{2 |x + 2|}
		      \text.
	      \]
	      Becsüljük felül a törtünket. Ehhez a nevezőt kell alulról becsülni,
	      ezért szorítkozzunk a $0$ egy környezetére: ha $|x| < 1$, akkor
	      $1 < x + 2 < 3$, így $|x + 2| > 1$. Ekkor:
	      \[
		      \frac{|x|}{2 |x + 2|}
		      < \frac{|x|}{2}
		      < \varepsilon
		      \qquad \Rightarrow \qquad
		      |x| < 2 \varepsilon
		      \text.
	      \]
	      A megoldás tehát:
	      \[
		      \delta(\varepsilon) = \min \left\{ 1; 2 \varepsilon \right\}
		      \text.
	      \]

	\item Az egyszerűsített alak az $x = -3$ pontban folytonos, így
	      behelyettesíthetünk:
	      \[
		      \lim_{x \to -3} f(x)
		      = \lim_{x \to -3} \frac{x + 3}{x + 2}
		      = \frac{0}{-1}
		      = 0
		      \text.
	      \]

	\item Osszuk el a számlálót és a nevezőt is $x$-szel:
	      \[
		      \lim_{x \to +\infty} f(x)
		      = \lim_{x \to +\infty} \frac{1 + \sfrac{3}{x}}{1 + \sfrac{2}{x}}
		      = 1
		      \text.
	      \]
\end{enumerate}
\end{document}
